\documentclass[UTF8,a4paper,12pt,fontset=none]{ctexart}

\usepackage{amsmath,amssymb,bm}
\usepackage{amsthm}
\usepackage{graphicx}
\usepackage{geometry}
\usepackage{booktabs}
\usepackage{float}
\usepackage{hyperref}
\usepackage[backend=biber,style=gb7714-2015,sorting=none]{biblatex}

\setCJKmainfont{Songti SC}
\setCJKsansfont{PingFang SC}
\setCJKmonofont{Songti SC}

\addbibresource{/Users/a1/sys/profiles/backup/zotero/我的文库.bib}
\addbibresource{refs/local.bib}

\geometry{left=2.6cm,right=2.6cm,top=2.6cm,bottom=2.6cm}
\hypersetup{colorlinks=true,linkcolor=blue,citecolor=blue,urlcolor=blue}
\setlength{\parindent}{2em}
\setlength{\parskip}{0.25em}

% --- Theorem environments ---
\theoremstyle{definition}
\newtheorem{assumption}{假设}[section]
\newtheorem{definition}{定义}[section]
\newtheorem{problem}{问题}[section]

\theoremstyle{plain}
\newtheorem{theorem}{定理}[section]
\newtheorem{lemma}{引理}[section]
\newtheorem{proposition}{命题}[section]
\newtheorem{corollary}{推论}[section]

\theoremstyle{remark}
\newtheorem{remark}{注}[section]

\title{基于光滑周期延迟反馈的自由漂浮空间机械臂\\预设时间轨迹跟踪控制}
\date{}

\begin{document}

\maketitle

\begin{abstract}
针对捕获前自由漂浮空间机械臂基座不受控、系统动力学强耦合且末端轨迹跟踪精度要求高的问题，本文提出一种基于光滑周期延迟反馈（smooth periodic delayed feedback, PDF）的关节空间轨迹跟踪控制方法。首先依据动量守恒关系，将含基座反作用运动的系统动力学约化为等效关节空间模型；针对有界集总扰动，设计预设时间扰动观测器生成 \(\hat d(t)\)，用于在指定观测时间后补偿匹配扰动；随后构造动力学补偿项，在误差动力学下建立具有可控规范形的辅助输入通道；最后引入周期时变延迟反馈，构造由当前误差和延迟误差共同组成的预设时间跟踪控制律。本文在定理--证明框架下给出以下理论结果：(i) 在连续时间、精确模型和精确补偿条件下，闭环误差可在预设时间 \(T_o+2h\) 后精确归零；(ii) 在存在观测误差、模型近似、执行器限制和数值离散时，结论退化为预设时间附近的实际有界跟踪，并给出明确的误差上界表达式。基于七自由度空间机械臂参数的名义等效关节空间模型 MATLAB 仿真验证了控制结构在有界扰动下的数值可行性。
\end{abstract}

\noindent\textbf{关键词：}自由漂浮空间机械臂；轨迹跟踪；预设时间控制；周期延迟反馈；扰动观测器

\section{引言}

随着在轨服务、空间碎片清理和失效航天器维护需求的增长，空间机械臂已成为执行非合作目标接近、捕获与操作任务的重要装备\cite{liu2021review}。与地面固定基座机械臂不同，捕获前空间机械臂通常处于自由漂浮状态，基座不直接受控，机械臂关节运动会通过系统动量守恒关系诱导基座线速度和角速度变化，进而改变末端执行器的实际运动。严宇新针对在轨捕获任务指出，空间机械臂动力学模型具有高度非线性和强耦合特性，并且轨迹跟踪过程中还会受到模型参数不确定、未知时变扰动和执行器故障等因素影响\cite{yan}。因此，如何在基座不直接施加控制力矩的条件下实现高精度、可预测收敛时间的轨迹跟踪，是自由漂浮空间机械臂控制中的关键问题。

自由漂浮空间机械臂的控制难点首先来自建模层面的非完整约束。在线动量和角动量初值为零的自由漂浮模式下，基座速度可由机械臂关节速度唯一诱导，末端速度不再由固定基座雅可比矩阵决定，而应由包含基座反作用运动的广义雅可比矩阵描述。

Umetani 和 Yoshida 提出的广义雅可比矩阵为自由漂浮空间机械臂的速度级建模和分解运动速度控制奠定了经典基础\cite{umetani1989}；Vafa 和 Dubowsky 的虚拟机械臂方法从动量守恒角度系统刻画了空间机械臂运动学与动力学耦合\cite{vafa1990}；Papadopoulos 和 Dubowsky 进一步揭示了自由漂浮空间机械臂中由动力学耦合引起的动态奇异问题，并系统总结了自由漂浮与自由飞行两类空间机器人的运动学与动力学特性\cite{papadopoulos1993,dubowsky1993}。在动力学建模与控制层面，Parlaktuna 和 Ozkan 提出了基于动力学等效机械臂的自适应计算力矩控制方案\cite{parlaktuna2004}；Xu 等进一步设计了自适应零反作用运动控制器，在跟踪末端轨迹的同时主动抑制基座姿态扰动\cite{xu2016}。严宇新论文中由动量守恒方程推导出基座--机械臂雅可比矩阵 \(J_{bm}\)，并进一步得到 \(J_g=J_bJ_{bm}+J_m\) 的自由漂浮运动学关系；在动力学层面，通过消去基座加速度，可将系统写成 \(H_g(q)\ddot q+C_g(q,\dot q)\dot q=\tau_m\) 的等效关节空间形式\cite{yan}。这一约化形式为把自由漂浮系统转化为关节空间控制问题提供了基础。

已有研究围绕轨迹规划、约束处理和鲁棒跟踪展开。针对自由漂浮模式下的运动学冗余、动力学奇异和关节速度约束，Wang 等采用粒子群优化生成满足约束的关节轨迹\cite{wang2015}；宁昕和武耀发针对传统控制方法难以同时处理自由漂浮空间机器人轨迹跟踪约束的问题，基于拉格朗日动力学模型构造扩展状态空间模型，并利用模型预测控制实现对末端期望位置和速度的跟踪\cite{ning2019}；针对捕获前反作用抑制和参数不确定问题，也有研究将反作用零空间规划、时延估计和自适应滑模控制结合，以降低对精确动力学模型的依赖\cite{yu2022}。Yao 基于障碍李雅普诺夫函数和径向基神经网络，为自由飞行空间机械臂设计了全状态和输出反馈自适应预设性能控制器，确保位置跟踪误差在参数不确定、外部扰动和输入饱和条件下始终保持在预定性能边界内\cite{yao2021}。这些方法在跟踪精度和鲁棒性方面各有侧重，但关于跟踪误差的收敛时间通常依赖于控制器参数与初始状态的复杂关系，难以在设计阶段独立指定。

近年来，固定时间（fixed-time）、预定义时间（predefined-time）和指定时间（prescribed-time）控制进一步受到关注，其共同优势在于可在设计阶段给出与初始误差弱相关或无关的收敛时间估计。在非线性系统的固定时间控制方面，Polyakov 提出了基于非连续齐次性的固定时间镇定框架\cite{polyakov2012}；Zuo 等进一步将固定时间方法推广至多智能体系统一致性控制\cite{zuo2019}。对于机器人系统，Su 等利用非奇异终端滑模实现了机械臂的固定时间轨迹跟踪控制\cite{su2020}；相关研究已将预定义时间稳定思想用于自由飞行空间机器人 SE(3) 位姿跟踪\cite{xuPredefinedTimeTrackingControl2024}，也将固定时间容错控制用于含执行器故障和性能约束的自由漂浮空间机械臂\cite{li2025}。然而，上述固定时间或预定义时间方法大多依赖非光滑终端滑模面或分数幂反馈项，在收敛时刻附近存在控制输入跳变或增益奇异问题，对工程实现构成挑战。

离散周期延迟反馈（delayed feedback control, DFC）概念最早由 Pyragas 提出，用于镇定混沌系统中的不稳定周期轨道\cite{pyragas1992}。其基本思想是通过系统当前状态与延迟一个周期长度的状态之差构造反馈信号，在不改变目标轨道的前提下实现镇定。延拓至连续时间框架后，Zhou、Michiels 和 Chen 提出的光滑周期延迟反馈（smooth periodic delayed feedback, PDF）方法表明，对于可控线性时滞系统，若预设收敛时间满足相应可解条件，可通过周期延迟反馈实现固定时间镇定；其反馈增益可选为连续、连续可微甚至光滑函数，从而避免传统指定时间高增益方法在终端时刻出现的增益奇异\cite{zhouFixedtimeStabilizationLinear2022}。该方法的核心不是构造非光滑终端滑模面，而是利用周期延迟系统的有限维单周期状态转移矩阵（又称单值矩阵），并通过使该矩阵零化或幂零来获得固定时间收敛。该方法在线性系统和线性化模型框架下具有理论简洁性和优越的时域性能，但其适用性需要被控系统首先转换为线性可控形式。

需要说明的是，固定时间收敛结论从非线性原系统到线性误差通道的迁移依赖于扰动补偿的精度。预设时间扰动观测器（prescribed-time disturbance observer, PTDO）是达成这一目的的重要工具。已有研究表明，通过时变增益构造可保证观测误差在预设时间内精确收敛至零\cite{jiangPrescribedtimeDisturbanceObserver2024}。在机器人控制领域，扰动观测器结合计算力矩补偿的方法已得到广泛验证，但将其与固定时间反馈结构显式衔接并给出联合预设时间收敛保证的相关工作尚不充分。

基于以上背景，本文采用"扰动观测--动力学补偿--误差线性化--周期延迟反馈"的结构：先依据动量守恒关系将捕获前自由漂浮空间机械臂约化为等效关节空间模型，再通过预设时间扰动观测器得到集总扰动估计 \(\hat d(t)\)，随后利用计算力矩补偿得到从辅助加速度到跟踪误差的线性可控通道，最后在该误差通道中嵌入光滑周期延迟反馈。这样既保留 PDF 方法的固定时间镇定机理，又使其作用对象与原理论假设保持一致。

本文主要工作如下：第一，基于自由漂浮空间机械臂的动量守恒约束，给出从基座反作用运动到广义雅可比矩阵、等效关节空间动力学的推导链条；第二，引入预设时间扰动观测补偿机制，说明 \(\hat d(t)\) 可由等效加速度扰动观测得到，并给出观测--PDF 两阶段设计时序；第三，构造面向 PDF 设计的误差通道规范化方法，将关节轨迹跟踪问题转化为线性可控误差系统镇定问题；第四，结合 PDF 理论中的 \(R_h(t)\)、可控 Gramian 和单周期状态转移矩阵零化条件，给出关节力矩控制律及其固定时间收敛条件；第五，基于七自由度空间机械臂参数搭建 MATLAB 仿真，验证所提控制结构的可运行性和有界扰动下的跟踪性能。

\section{自由漂浮空间机械臂建模}
\label{sec:modeling}

本节建立捕获前自由漂浮空间机械臂的等效关节空间动力学模型。在给出完整动力学方程后，利用动量守恒消去基座加速度，得到一种适用于后续控制器设计的降阶形式。

\subsection{符号与记法}

考虑一个由基座和 \(n\) 个旋转关节组成的自由漂浮空间机械臂。基座广义速度定义为
\[
\dot x_b=[v_b^\top,\ \omega_b^\top]^\top\in\mathbb R^6,
\]
其中 \(v_b\in\mathbb R^3\) 和 \(\omega_b\in\mathbb R^3\) 分别为基座线速度和角速度。关节变量记为 \(q=[q_1,\ldots,q_n]^\top\in\mathbb R^n\)。末端执行器广义速度为
\[
\dot x_e=[v_e^\top,\ \omega_e^\top]^\top\in\mathbb R^6.
\]

\subsection{完整动力学模型}

在无外力、无外力矩的自由漂浮条件下，系统动能可写为二次型形式
\begin{equation}
T=\frac12
\begin{bmatrix}\dot x_b^\top & \dot q^\top\end{bmatrix}
\begin{bmatrix}
H_b(q) & H_{bm}(q)\\
H_{bm}^\top(q) & H_m(q)
\end{bmatrix}
\begin{bmatrix}\dot x_b\\ \dot q\end{bmatrix},
\label{eq:kinetic_energy_ffsm}
\end{equation}
其中 \(H_b(q)\in\mathbb R^{6\times6}\)、\(H_m(q)\in\mathbb R^{n\times n}\) 和 \(H_{bm}(q)\in\mathbb R^{6\times n}\) 分别为基座惯性块、机械臂惯性块及基座--机械臂耦合惯性块。式~\eqref{eq:kinetic_energy_ffsm}~中的各惯性块在统一参考系下定义；若采用不同空间坐标约定，应首先通过相应坐标变换统一动量表达。

基于拉格朗日力学，完整动力学方程可写为分块形式
\begin{equation}
\begin{bmatrix}
H_b & H_{bm}\\
H_{bm}^\top & H_m
\end{bmatrix}
\begin{bmatrix}\ddot x_b\\ \ddot q\end{bmatrix}
+
\begin{bmatrix}
C_b & C_{bm}\\
C_{bm}^\top & C_m
\end{bmatrix}
\begin{bmatrix}\dot x_b\\ \dot q\end{bmatrix}
=
\begin{bmatrix}0\\ \tau\end{bmatrix},
\label{eq:full_ffsm_dynamics}
\end{equation}
其中 \(\tau=[\tau_1,\ldots,\tau_n]^\top\in\mathbb R^n\) 为关节驱动力矩。式~\eqref{eq:full_ffsm_dynamics}~的第一行体现了自由漂浮的物理条件：基座不直接施加控制力和控制力矩。\(C\) 矩阵的各子块由 Christoffel 符号或递推 Newton--Euler 方法确定，满足 \(\dot H-2C\) 的斜对称性质（在适当定义下）。

\subsection{动量守恒与广义雅可比关系}

\noindent\textbf{假设1（零初始动量）：} 系统初始总线动量和总角动量均为零。\hfill\(\blacksquare\)

在假设~1 下，动量守恒约束可写为
\begin{equation}
H_b(q)\dot x_b + H_{bm}(q)\dot q = 0.
\label{eq:momentum_constraint}
\end{equation}

\noindent\textbf{假设2（基座惯性可逆）：} 在所考虑的工作区域内，\(H_b(q)\) 一致正定。\hfill\(\blacksquare\)

在假设~2 下，由式~\eqref{eq:momentum_constraint}~可得基座速度与关节速度之间的线性映射关系
\begin{equation}
\dot x_b = -H_b^{-1}(q)H_{bm}(q)\dot q \triangleq J_{bm}(q)\dot q,
\label{eq:base_joint_jacobian}
\end{equation}
其中 \(J_{bm}(q)\in\mathbb R^{6\times n}\) 称为基座--机械臂雅可比矩阵。式~\eqref{eq:base_joint_jacobian}~表明，在自由漂浮模式下，基座速度不是独立控制量，而是由关节运动通过动量守恒诱导产生的反作用运动。

设末端速度由运动学链关系给出
\begin{equation}
\dot x_e = J_b(q)\dot x_b + J_m(q)\dot q,
\label{eq:end_velocity_common}
\end{equation}
其中 \(J_b(q)\in\mathbb R^{6\times6}\) 和 \(J_m(q)\in\mathbb R^{6\times n}\) 分别为相对于基座坐标的末端雅可比分量。将式~\eqref{eq:base_joint_jacobian}~代入式~\eqref{eq:end_velocity_common}，得到自由漂浮空间机械臂的广义雅可比关系
\begin{equation}
\dot x_e = \bigl[J_m(q) + J_b(q)J_{bm}(q)\bigr]\dot q \triangleq J_g(q)\dot q.
\label{eq:generalized_jacobian}
\end{equation}
该关系表明末端运动同时受关节运动和基座反作用运动的共同影响。当 \(J_g(q)\) 接近秩亏时，关节速度放大、基座反作用增强以及末端可操作性退化可能同时出现。

\subsection{等效关节空间动力学}

将式~\eqref{eq:momentum_constraint}~对时间求导可得
\begin{equation}
H_b\ddot x_b + H_{bm}\ddot q + \dot H_b\dot x_b + \dot H_{bm}\dot q = 0.
\label{eq:momentum_derivative}
\end{equation}
消去 \(\ddot x_b\)：由式~\eqref{eq:full_ffsm_dynamics}~的第一行解出 \(\ddot x_b\) 的表达式，代入式~\eqref{eq:momentum_derivative}~并与式~\eqref{eq:full_ffsm_dynamics}~的第二行联立，经过代数运算（该运算涉及各惯性块及其时间导数和 Coriolis 项的重新组合），得到等效关节空间动力学
\begin{equation}
M_e(q)\ddot q + C_e(q,\dot q)\dot q = \tau + d(t).
\label{eq:joint_model}
\end{equation}
其中等效惯性矩阵为 Schur 补
\begin{equation}
M_e(q) = H_m(q) - H_{bm}^\top(q)H_b^{-1}(q)H_{bm}(q).
\label{eq:generalized_inertia}
\end{equation}
\(C_e(q,\dot q)\dot q \in\mathbb R^n\) 表示由完整递推动力学、动量约束求导以及坐标变换共同确定的等效速度相关项；该项的具体矩阵分解并不唯一，但应使式~\eqref{eq:joint_model}~与完整动力学方程在零初始动量条件下一致。\(d(t)\in\mathbb R^n\) 为集总扰动项，汇总参数摄动、未建模动力学、外部扰动以及执行器误差等所有非理想因素。

\noindent\textbf{假设3（等效惯性正定性）：} 在所考虑的工作区域内，\(M_e(q)\) 一致正定且可逆。\hfill\(\blacksquare\)

假设~3 保证式~\eqref{eq:joint_model}~对控制器设计是良定义的。由于式~\eqref{eq:generalized_inertia}~是 Schur 补形式，当完整惯性矩阵正定时该假设自然成立。

为与后续控制器记号一致，将关节变量统一记为 \(q_m\in\mathbb R^n\)，将等效动力学写为
\begin{equation}
M_e(q_m)\ddot q_m + C_e(q_m,\dot q_m)\dot q_m = \tau + d(t).
\label{eq:joint_model_relabelled}
\end{equation}
式~\eqref{eq:joint_model_relabelled}~即为本文后续控制律设计的基础模型。该约化模型保留了基座反作用对关节动力学的等效影响，将原 \(6+n\) 维耦合系统降为 \(n\) 维关节空间形式，为计算力矩补偿和误差通道的线性化提供了分析基础。

\begin{remark}[空间 CRBA 数值实现]
  在实际数值实现中，完整惯性矩阵 \(H_{\text{full}}(q)\in\mathbb R^{(6+n)\times(6+n)}\) 可通过空间复合刚体算法（Composite Rigid Body Algorithm, CRBA）计算：将各连杆的空间惯性矩阵 \(I_i\in\mathbb R^{6\times6}\) 经伴随变换表达在基座坐标系中求和，即
  \[
  H_{\text{full}} = \sum_{i=0}^{n} J_i^\top I_i J_i,
  \]
  其中 \(J_i\in\mathbb R^{6\times(6+n)}\) 为连杆 \(i\) 在自身坐标系中的空间雅可比。该方法的计算复杂度为 \(O(n^2)\)，适用于在线实现。科里奥利/离心项 \(C_{\text{full}}\dot q_{\text{full}}\) 可通过 Christoffel 符号的数值微分计算。完整的前向动力学求解步骤为
  \[
  \ddot q_{\text{full}} = H_{\text{full}}^{-1}\bigl([0_{6\times1}^\top,\;\tau^\top]^\top - C_{\text{full}}\dot q_{\text{full}}\bigr),
  \]
  其中前 \(6\) 个加速度分量对应基座自由度，后 \(n\) 个对应关节自由度。提取 Schur 补 \(M_e = H_m - H_{bm}^\top H_b^{-1} H_{bm}\) 及等效科里奥利向量 \(C_e\dot q = C_m - H_{bm}^\top H_b^{-1} C_b\) 即可用于控制器设计。本文第~\ref{sec:simulation}~节的仿真采用基于该方法参数的正定名义关节空间模型，完整的 CRBA 递推动力学实现留待后续工作验证。
\end{remark}
\section{问题描述}

给定足够光滑且有界的关节期望轨迹 \(q_d(t),\dot q_d(t),\ddot q_d(t)\)。定义位置误差、速度误差和误差状态为
\[
e=q_m-q_d,\qquad \dot e=\dot q_m-\dot q_d,
\qquad z=\begin{bmatrix}e^\top&\dot e^\top\end{bmatrix}^\top\in\mathbb R^{2n}.
\]
本文首先研究关节空间轨迹跟踪问题，控制目标是设计关节力矩 \(\tau\)，使误差在给定时间上界后满足
\begin{equation}
e(t)=0,\qquad \dot e(t)=0,\qquad t\ge T.
\label{eq:tracking_goal}
\end{equation}
其中，\(T\) 由观测阶段和 PDF 阶段的设计参数共同确定。需要区分的是，在理想连续时间、精确模型和精确扰动补偿条件下，式 \eqref{eq:tracking_goal} 表示预设时间精确跟踪；当存在模型误差、持续未补偿扰动、执行器饱和或数值离散时，本文只主张预设时间附近的实际有界跟踪或小邻域收敛。

本文采用如下假设。

\textbf{假设1：} 在所考虑工作区域内，\(M_e(q_m)\) 一致正定且可逆，\(M_e(q_m)\) 与 \(C_e(q_m,\dot q_m)\dot q_m\) 可由完整模型获得，或由名义模型近似获得。

\textbf{假设2：} 参考轨迹 \(q_d,\dot q_d,\ddot q_d\) 有界且连续，关节位置和速度可测，且参考信号在控制器启动时具有一致的历史定义。

\textbf{假设3：} 人工延迟项 \(z(t-h)\) 可由历史缓存获得，其中 \(h>0\) 为设计延迟。延迟初值函数在 \([-h,0]\) 上有界且分段连续。

\textbf{假设4：} 对于预设时间扰动观测器，等效加速度扰动
\[
\Delta_a(t)=M_e^{-1}(q_m)d(t)
\]
及其一阶导数有界；其上界可以未知。

\textbf{假设5：} 严格的预设时间精确跟踪结论只在连续时间实现、无执行器饱和、模型补偿精确以及观测器按理论条件实现时成立。存在观测误差、模型不确定性、采样和量化误差时，相关项统一视为残余扰动进行实际有界性分析。

上述假设明确了本文理论结果的适用范围。后文先在理想条件下建立 PDF 闭环的预设时间精确跟踪结论，再讨论残余扰动对实际性能的影响。末端任务空间轨迹若由广义雅可比映射生成关节参考，则属于关节空间控制器的上层参考生成问题，本文的主要理论对象仍为式 \eqref{eq:tracking_goal} 所示的关节空间误差系统。

\section{控制算法设计}

\subsection{预设时间扰动观测器设计}

对等效关节空间模型 \eqref{eq:joint_model}，先设计扰动观测器来生成后续动力学补偿所需的 \(\hat d(t)\)。将关节速度记为
\[
\chi=\dot q_m .
\]
由 \eqref{eq:joint_model} 可得
\begin{equation}
\dot\chi=u_o+\Delta_a,
\qquad
u_o=M_e^{-1}(q_m)\bigl[\tau-C_e(q_m,\dot q_m)\dot q_m\bigr],
\qquad
\Delta_a=M_e^{-1}(q_m)d(t).
\label{eq:observer_channel}
\end{equation}
其中 \(\Delta_a\) 为匹配到加速度通道的集总扰动。该形式与预设时间扰动观测器的标准受扰单积分器模型一致\cite{jiangPrescribedtimeDisturbanceObserver2024}。若能够在预设时间 \(T_o\) 内得到 \(\hat\Delta_a\)，则力矩扰动估计可由
\begin{equation}
\hat d(t)=M_e(q_m)\hat\Delta_a(t)
\label{eq:dhat_from_observer}
\end{equation}
给出。

设 \(z_1,z_2\in\mathbb R^n\) 分别为 \(\chi\) 和 \(\Delta_a\) 的估计值，取 \(\hat\Delta_a=z_2\)。定义可测复合误差
\[
\varepsilon_1=\chi-z_1-\bar\xi(t),
\]
其中 \(\bar\xi(t)\in\mathbb R^n\) 为调节函数，满足有界光滑、单调衰减且 \(t\ge T_o\) 时 \(\bar\xi(t)=0\)。例如可取
\[
\bar\xi(t)=
\begin{cases}
\bar\xi_0(T_o-t)^2, & 0\le t<T_o,\\
0, & t\ge T_o.
\end{cases}
\]
记 \(\operatorname{sig}^{\alpha}(x)=\operatorname{sgn}(x)|x|^\alpha\)，并按元素作用于向量。参考文献\cite{jiangPrescribedtimeDisturbanceObserver2024}，构造如下预设时间扰动观测器：
\begin{equation}
\left\{
\begin{aligned}
\dot z_1
&=z_2+\frac{\pi}{\eta T_c}
\phi_1\left(\frac{\varepsilon_1}{\sigma}\right)
-\dot{\bar\xi}(t)+\bar\xi(t)+u_o,\\
\dot z_2
&=\frac{\pi}{\sigma\eta T_c}
\phi_2\left(\frac{\varepsilon_1}{\sigma}\right)
-\dot{\bar\xi}(t),
\end{aligned}
\right.
\label{eq:ptdo}
\end{equation}
其中 \(T_c\le T_o\)、\(\eta\in(0,1)\)、\(\sigma>0\)，且
\[
\begin{aligned}
\phi_1(x)&=\operatorname{sig}^{1-\eta/2}(x)+\operatorname{sig}^{1+\eta/2}(x),\\
\phi_2(x)&=\operatorname{sig}^{2-\eta}(x)+\operatorname{sig}^{2+\eta}(x)+\sigma\operatorname{sgn}(x).
\end{aligned}
\]
在 \(\Delta_a\) 与 \(\dot\Delta_a\) 有界、连续时间实现且观测器满足相应理论条件时，观测误差可在设计时间 \(T_o\) 内收敛到零。该结论仅适用于连续时间理想观测器；在采样、量化、测量噪声和执行器限幅存在时，观测误差通常只能收敛到与采样步长、数值积分误差和未建模动态有关的小邻域。参数 \(T_o\) 和 \(T_c\) 不宜选得过小，否则可能带来较大的观测峰值和控制输入峰值。

\subsection{动力学补偿与误差通道规范化}

在获得 \(\hat d(t)\) 后，采用如下动力学补偿结构
\begin{equation}
\tau=M_e(q_m)v+C_e(q_m,\dot q_m)\dot q_m-\hat d(t),
\label{eq:computed_torque}
\end{equation}
其中 \(v\) 为辅助加速度输入。将式 \eqref{eq:computed_torque} 代入式 \eqref{eq:joint_model} 可得
\[
M_e(q_m)(\ddot q_m-v)=d(t)-\hat d(t).
\]
当观测器已收敛并满足 \(\hat d(t)=d(t)\) 时，且 \(M_e(q_m)\) 正定可逆，因此
\(
\ddot q_m=v.
\)

令
\(
v=\ddot q_d+\nu ,
\)
并定义
\(
e=q_m-q_d\),
\(z=[e, \dot e]^T.
\)
则
\(
\ddot e=\ddot q_m-\ddot q_d=\nu,
\)
误差系统可写成二阶积分链形式
\begin{equation}
\dot z=Az+B\nu ,
\label{eq:linear_error}
\end{equation}
其中
\[
A=\begin{bmatrix}0&I_n\\0&0\end{bmatrix},\qquad
B=\begin{bmatrix}0\\I_n\end{bmatrix}.
\]
矩阵对 \((A,B)\) 可控，因此 PDF 增益的设计对象由原非线性动力学转移为线性可控误差通道。这个步骤是本文方法成立的关键：PDF 理论只需作用在 \eqref{eq:linear_error} 上，而不直接处理自由漂浮机械臂的非线性耦合项。

\subsection{PDF 线性误差系统设计}

文献\cite{zhouFixedtimeStabilizationLinear2022}考虑的基本对象是可控线性时滞系统，其固定时间镇定依赖如下周期延迟结构：
\begin{equation}
\dot x(t)=F(t)x(t)+G(t)x(t-h),
\label{eq:periodic_delay_general}
\end{equation}
其中 \(F(t),G(t)\) 为 \(2h\) 周期矩阵，且 \(G(t)=0\) 在每个周期的前半段 \([0,h]\) 成立。由分步法可知，系统在一个周期后的状态满足
\(
x(2(k+1)h)=\Delta(h)x(2kh),
\)
其中单周期状态转移矩阵（又称单值矩阵）为
\begin{equation}
\Delta(h)=\Phi(2h,0)+\int_h^{2h}\Phi(2h,s)G_0(s)\Phi(s-h,0)\,ds.
\label{eq:period_map}
\end{equation}
若 \(\Delta(h)\) 为幂零矩阵，即存在最小整数 \(\nu_0\ge1\) 使 \(\Delta^{\nu_0}(h)=0\)，则 \(x(t)=0,\ t\ge 2\nu_0h\)。特别地，若通过增益设计使 \(\Delta(h)=0\)，则系统可在一个完整周期后归零\cite{zhouFixedtimeStabilizationLinear2022}。

为将该思想用于式 \eqref{eq:linear_error}，先取常值反馈 \(K_0\) 使
\(
A_c=A+BK_0
\)
成为期望的基准闭环矩阵，再引入周期延迟反馈
\begin{equation}
\nu(t)=K_0z(t)-K_c(t)z(t-h).
\label{eq:pdf_error_input}
\end{equation}
此时误差闭环为
\begin{equation}
\dot z(t)=A_cz(t)-BK_c(t)z(t-h).
\label{eq:closed_delay_error}
\end{equation}
令 \(R_h(t)\) 为满足
\[
R_h(t+2h)=R_h(t)
\]
的 \(2h\) 周期光滑函数，并满足
\[
R_h(t)=0,\quad t\in[0,h],\qquad
R_h(t)\ge0,\quad t\in[h,2h],
\]
且在 \([h,2h]\) 的某个具有非零测度的子区间上严格为正。本文采用标量函数情形。根据文献\cite{zhouFixedtimeStabilizationLinear2022}的构造，可定义
\begin{equation}
W_c(A_c,h)=\int_h^{2h}e^{-A_cs}BR_h(s)B^\top e^{-A_c^\top s}\,ds,
\label{eq:wc_matrix}
\end{equation}
并取
\begin{equation}
K_c(t)=R_h(t)B^\top e^{-A_c^\top t}W_c^{-1}(A_c,h)e^{A_c(h-t)}.
\label{eq:kc_pdf}
\end{equation}
由于 \((A_c,B)\) 与 \((A,B)\) 同可控，且 \(R_h(t)\) 在后半周期的一个非零测度子区间上严格为正，连续时间条件下 \(W_c(A_c,h)\) 为正定并可逆。将 \eqref{eq:kc_pdf} 代入 \eqref{eq:period_map} 可使闭环的单周期状态转移矩阵满足 \(\Delta(h)=0\)，从而得到单周期零化性质。该结论依赖精确的 Gramian 计算和连续时间延迟实现；数值积分、矩阵病态和采样误差会使实际闭环仅近似满足零化条件。式 \eqref{eq:kc_pdf} 也说明 PDF 增益可通过离线积分计算，在线控制只需读取当前误差和历史误差。

结合式 \eqref{eq:computed_torque} 和式 \eqref{eq:pdf_error_input}，PDF 关节力矩控制律可写为
\begin{equation}
\tau
=
M_e(q_m)
\left[
\ddot q_d
+
K_0z(t)
-
K_c(t)z(t-h)
\right]
+
C_e(q_m,\dot q_m)\dot q_m
-
\hat d(t).
\label{eq:final_controller}
\end{equation}

\subsection{固定时间抗饱和补偿器}

实际应用中，关节执行器存在力矩饱和限制，即
\begin{equation}
\tau_i = \operatorname{sat}(\tau_{ci}) =
\begin{cases}
\tau_{i,\max}\,\operatorname{sgn}(\tau_{ci}), & |\tau_{ci}|>\tau_{i,\max},\\
\tau_{ci}, & |\tau_{ci}|\le\tau_{i,\max},
\end{cases}
\label{eq:saturation_model}
\end{equation}
其中 \(\tau_{ci}\) 为第 \(i\) 个关节的计算控制力矩，\(\tau_{i,\max}>0\) 为力矩幅值上限。饱和将导致实际施加力矩与计算力矩之间产生偏差 \(\Delta\tau_i=\tau_i-\tau_{ci}\)。映射到加速度通道的归一化饱和误差为
\begin{equation}
\Delta_u(t)=M_e^{-1}(q_m)\Delta\tau(t).
\label{eq:norm_saturation_error}
\end{equation}

为抑制输入饱和对跟踪性能的影响，参考文献\cite{dou2025fractional}，引入如下固定时间抗饱和补偿器：
\begin{equation}
\dot\eta = -L_\eta\frac{k_{\eta1}}{\rho_1}
\tanh\bigl(|\eta|^{1-\rho_1}\bigr)
\cosh\bigl(|\eta|^{2}\bigr)
\operatorname{sgn}(\eta)
- L_\eta k_{\eta2}
\operatorname{sig}^{\rho_2}(\eta)
+ \Delta_u,
\label{eq:as_dynamics}
\end{equation}
其中 \(\eta\in\mathbb R^n\) 为辅助变量，\(0<\rho_1,\rho_2<1\)，\(k_{\eta1},k_{\eta2}>0\) 为增益系数，\(L_\eta>0\) 为自适应增益。\(\operatorname{sig}^{\alpha}(x)=\operatorname{sgn}(x)|x|^{\alpha}\) 按元素作用于向量。该补偿器的核心性质由以下定理给出：

\textbf{定理1：} 当 \(\Delta_u=0\) 时，辅助变量 \(\eta\) 将在固定时间内收敛至零，其上界为
\begin{equation}
T_{\eta,\max} < 
\frac{\rho_1 I}{k_{\eta1}\tanh(1)} + 
\frac{1}{k_{\eta2}(1-\rho_2)},
\end{equation}
其中 \(I=\int_1^{\infty}\frac{dz}{\cosh(z^2)}<\infty\) 为收敛常数。

\textbf{证明：} 令 \(z=|\eta|\)，由式 \eqref{eq:as_dynamics} 可得
\[
\dot z = -L_\eta\frac{k_{\eta1}}{\rho_1}
\tanh(z^{1-\rho_1})
\cosh(z^2)
- L_\eta k_{\eta2} z^{\rho_2}.
\]
当 \(z\ge 1\) 时，\(\tanh(z^{1-\rho_1})\ge\tanh(1)\) 且 \(\cosh(z^2)\ge e^{z^2}/2\)，因此
\[
T_3 = \int_1^{z_0}\frac{dz}{|\dot z|}
\le \frac{\rho_1}{L_\eta k_{\eta1}\tanh(1)}
\int_1^{\infty}\frac{dz}{\cosh(z^2)}
= \frac{\rho_1 I}{L_\eta k_{\eta1}\tanh(1)}.
\]
当 \(z\le 1\) 时，\(\dot z\le -L_\eta k_{\eta2}z^{\rho_2}\)，解得
\(
T_4 < 1/(L_\eta k_{\eta2}(1-\rho_2))
\)。
因此总收敛时间上界为 \(T_{\eta,\max}=T_3+T_4\)，与初始值 \(z_0\) 无关。\hfill\(\square\)

将抗饱和补偿器与 PDF 控制器结合，得到如下修正控制律：
\begin{equation}
\begin{aligned}
\tau_c &=
M_e(q_m)\bigl[
\ddot q_d + K_0z(t) - K_c(t)z(t-h) + \eta
\bigr] \\
&\qquad + C_e(q_m,\dot q_m)\dot q_m - \hat d(t).
\label{eq:final_controller_as}
\end{aligned}
\end{equation}
实际施加力矩为 \(\tau=\operatorname{sat}(\tau_c)\)。与原始控制律 \eqref{eq:final_controller} 相比，式 \eqref{eq:final_controller_as} 在辅助加速度中增加了 \(\eta\) 项。当 \(|M_e^{-1}\tau_c|>\tau_{\max}\) 时，饱和误差 \(\Delta_u\) 驱动 \(\eta\) 沿式 \eqref{eq:as_dynamics} 演化；一旦饱和解除，\(\eta\) 在固定时间 \(T_{\eta,\max}\) 内衰减至零，控制律退化为原始的 PDF 形式。

需要注意，若观测器与 PDF 控制器同时从 \(t=0\) 启动，则 \(0\le t<T_o\) 内的观测误差会作为残余扰动进入误差系统，严格 \(2h\) 零化结论不再直接成立。为保持理论闭环清晰，可采用两阶段时序：先在 \(T_o\) 内完成扰动观测，随后以局部时间 \(\vartheta=t-T_o\) 启动 PDF 增益 \(K_c(\vartheta)\)。此时理想连续时间条件下的总预设收敛时间上界为
\[
T_{\mathrm{tot}}=T_o+2h .
\]
在非理想实现中，该时间应理解为进入理论实际有界性分析的起始时间尺度，而非严格的数值归零时刻。

% \subsection{实现说明与鲁棒项}

% 在数值实现中，本文默认按照式 \eqref{eq:wc_matrix} 和式 \eqref{eq:kc_pdf} 离线计算 \(W_c(A_c,h)\) 与 \(K_c(t)\)，并在每个采样时刻按 \(2h\) 周期读取局部时间对应的 PDF 增益。为便于调参和消融对比，程序中还保留工程化平滑延迟项
% \[
% K_0=[-K_p\ -K_d],\qquad
% -K_c(t)=\rho R_h(t)[-K_{ph}\ -K_{dh}],
% \]
% 其中 \(K_p,K_d,K_{ph},K_{dh}\) 为正定对角矩阵，\(\rho>0\) 为延迟反馈幅值系数。该工程化形式仅保留 PDF 的等待--作用周期结构，并不等价于严格的单周期状态转移矩阵零化增益；论文中的理论固定时间结论仍以式 \eqref{eq:kc_pdf} 的 Gramian 型 PDF 增益为前提。

% 若存在观测误差或未建模高频扰动，也可进一步引入饱和型鲁棒补偿
% \[
% \tau_r=-k_d\,\operatorname{sat}(s/\varphi),
% \qquad
% s=\dot e+\Lambda e,
% \]
% 并将 \(\tau_r\) 叠加到 \eqref{eq:final_controller}。其中 \(k_d>0\)、\(\varphi>0\)、\(\Lambda=\Lambda^\top>0\)。需要指出，加入非线性鲁棒补偿后，闭环不再是文献\cite{zhouFixedtimeStabilizationLinear2022}中的纯线性延迟系统，严格固定时间精确归零结论需要单独证明；工程上通常得到对扰动鲁棒的实际固定时间有界结果。

\section{稳定性与性能分析}

本节在理想连续时间条件下给出闭环预设时间跟踪结论，并进一步说明模型误差、补偿误差和数值实现对该结论的影响。除特别说明外，假设观测器在 \(T_o\) 内完成收敛，且 PDF 控制器从 \(t=T_o\) 起以局部时间 \(\vartheta=t-T_o\) 启动。

\textbf{命题1：} 若预设时间扰动观测器满足其连续时间收敛条件，且 \(\Delta_a\) 与 \(\dot\Delta_a\) 有界，则存在设计时间 \(T_o>0\)，使得
\[
\hat\Delta_a(t)=\Delta_a(t),\qquad \hat d(t)=d(t),
\qquad \forall t\ge T_o.
\]

\textbf{说明：} 命题1是连续时间理想观测器结论。由
\[
\hat d=M_e(q_m)\hat\Delta_a,
\qquad
\Delta_a=M_e^{-1}(q_m)d
\]
可知，当 \(\hat\Delta_a=\Delta_a\) 时有 \(\hat d=d\)。在采样实现、测量噪声或观测器公式未完全满足理论条件时，不能直接将该结论解释为数值观测误差严格等于零。

\textbf{命题2：} 假设命题1成立。若 \((A,B)\) 可控，\(K_0\) 使
\[
A_c=A+BK_0
\]
为期望的基准闭环矩阵，并且周期增益 \(K_c(t)\) 按式 \eqref{eq:kc_pdf} 构造，使相应单周期状态转移矩阵满足
\[
\Delta_K(h)=0,
\]
则理想模型条件下，关节跟踪误差满足
\[
e(t)=0,\qquad \dot e(t)=0,
\qquad \forall t\ge T_o+2h.
\]

\textbf{证明：} 当 \(t\ge T_o\) 时，由命题1有 \(\hat d=d\)。将控制律 \eqref{eq:computed_torque} 代入式 \eqref{eq:joint_model}，得到 \(\ddot q_m=v\)。再令
\[
v=\ddot q_d+\nu,
\]
则有 \(\ddot e=\nu\)，并且
\[
\dot z=Az+B\nu,
\qquad
A=\begin{bmatrix}0&I_n\\0&0\end{bmatrix},
\quad
B=\begin{bmatrix}0\\I_n\end{bmatrix}.
\]
代入 \(\nu=K_0z-K_c(t)z(t-h)\)，得
\[
\dot z=A_cz-BK_c(t)z(t-h).
\]
在 PDF 周期的前半段 \(K_c(t)=0\)，因此可按分步法求得一个完整周期的状态映射。依据式 \eqref{eq:period_map} 以及式 \eqref{eq:kc_pdf} 的 Gramian 构造，闭环单周期状态转移矩阵满足
\[
\Delta_K(h)=0.
\]
因此局部时间达到 \(2h\) 后有 \(z(t)=0\)。由于局部时间为 \(\vartheta=t-T_o\)，故 \(t\ge T_o+2h\) 时 \(e(t)=0\) 且 \(\dot e(t)=0\)。\hfill\(\square\)

上述命题给出本文控制器的理想预设时间上界
\begin{equation}
T_{\mathrm{pre}}=T_o+2h.
\label{eq:prescribed_total_time}
\end{equation}
这里的“预设”体现在 \(T_o\) 和 \(h\) 可在控制器设计阶段选定，而不是由初始误差大小决定。需要注意，式 \eqref{eq:prescribed_total_time} 是连续时间理想结论；采样、力矩饱和和模型误差会使实际误差不再严格于该时刻归零。

\textbf{命题3：} 若采用一般周期延迟增益，并且相应单周期状态转移矩阵 \(\Delta(h)\) 为幂零矩阵，即存在 \(\nu_0\ge1\) 使
\[
\Delta^{\nu_0}(h)=0,
\]
则理想补偿条件下有
\[
e(t)=0,\qquad \dot e(t)=0,
\qquad \forall t\ge T_o+2\nu_0h.
\]
该结论来自周期状态映射
\[
z(2(k+1)h)=\Delta(h)z(2kh).
\]
式 \eqref{eq:kc_pdf} 对应更强的单周期零化情形，即 \(\nu_0=1\)。

\subsection{残余扰动下的实际跟踪性能}

当模型存在近似误差、扰动估计不完全或执行器受到饱和限制时，定义等效残余加速度扰动
\begin{equation}
\Delta_r(t)=M_e^{-1}(q_m)\bigl[d(t)-\hat d(t)+\Delta_m(t)+\Delta_\tau(t)\bigr],
\label{eq:residual_acceleration}
\end{equation}
其中 \(\Delta_m\) 表示名义模型与真实模型之间的等效动力学误差，\(\Delta_\tau\) 表示力矩指令经饱和或量化后产生的输入误差。此时误差系统可写为
\begin{equation}
\dot z(t)=A_cz(t)-BK_c(t)z(t-h)+B\Delta_r(t).
\label{eq:perturbed_delay_error}
\end{equation}
若 \(\Delta_r(t)\) 有界，则理想单周期零化映射被有界输入项取代，严格精确归零一般不再成立。下面给出一个基于单周期输入--状态映射的实际有界性结论。

定义 \(\mathcal S_h(\theta,s)\) 为齐次延迟系统
\[
\dot z(t)=A_cz(t)-BK_c(t)z(t-h)
\]
在一个 PDF 周期内从时刻 \(s\) 到时刻 \(\theta\) 的因果解算子，并定义周期内的扰动放大系数
\begin{equation}
\Gamma_h=\max_{0\le \theta\le 2h}
\int_0^\theta
\left\|\mathcal S_h(\theta,s)B\right\|\,\mathrm ds .
\label{eq:pdf_disturbance_gain}
\end{equation}
这里的范数可取诱导二范数。对于给定的 \(A_c\)、\(K_c\) 和 \(h\)，\(\Gamma_h\) 可通过离线数值积分或对相应延迟系统进行离散化计算。

\textbf{命题4：} 假设 PDF 增益满足理想单周期零化条件 \(\Delta_K(h)=0\)，且残余扰动满足
\[
\|\Delta_r\|_\infty
=\sup_{t\ge T_o}\|\Delta_r(t)\|\le\bar\Delta_r.
\]
则在连续时间实现下，完整 PDF 周期边界处的闭环误差满足
\begin{equation}
\|z(T_o+2(k+1)h)\|\le \Gamma_h\bar\Delta_r,
\qquad k=0,1,2,\ldots .
\label{eq:pdf_practical_bound}
\end{equation}
特别地，当 \(\Delta_r(t)\equiv0\) 时，式 \eqref{eq:pdf_practical_bound} 退化为理想的单周期精确归零结论。

\textbf{证明：} 令 \(z_k=z(T_o+2kh)\)，并将第 \(k\) 个 PDF 周期内的局部时间记为 \(\vartheta\in[0,2h]\)。变分公式给出
\[
z_{k+1}=\Delta_K(h)z_k+
\int_0^{2h}\mathcal S_h(2h,s)B\Delta_r(T_o+2kh+s)\,\mathrm ds.
\]
由 \(\Delta_K(h)=0\) 以及 \(\|\Delta_r\|_\infty\le\bar\Delta_r\)，有
\[
\|z_{k+1}\|
\le \int_0^{2h}
\|\mathcal S_h(2h,s)B\|\,\|\Delta_r(T_o+2kh+s)\|\,\mathrm ds
\le \Gamma_h\bar\Delta_r.
\]
该不等式对每个 \(k\ge0\) 成立，故得到式 \eqref{eq:pdf_practical_bound}。注意，该结论首先约束的是完整 PDF 周期的边界状态；若需要给出周期内任意时刻的误差界，还需额外定义周期内状态传播增益。\hfill\(\square\)

命题4给出了方向一所需的定量结论：完整 PDF 周期边界处的实际误差邻域大小由残余扰动上界 \(\bar\Delta_r\) 和 PDF 周期扰动增益 \(\Gamma_h\) 共同决定。若需要将结果推广到周期内任意时刻，还需补充周期内状态传播增益。进一步地，若数值实现或正则化计算只能保证
\[
\|\Delta_K(h)\|\le\varepsilon_\Delta<1,
\]
则在 PDF 周期边界处有递推不等式
\begin{equation}
\|z_{k+1}\|
\le \varepsilon_\Delta\|z_k\|
+\Gamma_h\bar\Delta_r,
\label{eq:pdf_approximate_recursion}
\end{equation}
从而
\begin{equation}
\limsup_{k\to\infty}\|z_k\|
\le \frac{\Gamma_h\bar\Delta_r}{1-\varepsilon_\Delta}.
\label{eq:pdf_ultimate_bound}
\end{equation}
式 \eqref{eq:pdf_ultimate_bound} 表明，Gramian 数值误差、采样误差和正则化误差会通过 \(\varepsilon_\Delta\) 放大残余扰动引起的稳态误差。由 \eqref{eq:residual_acceleration}，若
\[
\|M_e^{-1}(q_m)\|\le\bar m^{-1},\quad
\|d-\hat d\|\le\bar d_o,\quad
\|\Delta_m\|\le\bar\Delta_m,\quad
\|\Delta_\tau\|\le\bar\Delta_\tau,
\]
则可以采用保守估计
\begin{equation}
\bar\Delta_r\le\bar m^{-1}
\left(\bar d_o+\bar\Delta_m+\bar\Delta_\tau\right).
\label{eq:residual_bound_explicit}
\end{equation}
因此，理论误差界可进一步写成
\[
\|z(t)\|
\le \Gamma_h\bar m^{-1}
\left(\bar d_o+\bar\Delta_m+\bar\Delta_\tau\right)
\]
的形式；对于近似 PDF，则将右端替换为式 \eqref{eq:pdf_ultimate_bound} 给出的上界。

\subsection{抗饱和补偿器的稳定性分析}

当采用抗饱和补偿器 \eqref{eq:as_dynamics} 和修正控制律 \eqref{eq:final_controller_as} 时，式 \eqref{eq:residual_acceleration} 中的力矩输入误差 \(\Delta_\tau\) 由辅助变量 \(\eta\) 的动态补偿。此时残余扰动可重新写为
\begin{equation}
\Delta_r^{\mathrm{as}}(t)=M_e^{-1}(q_m)\bigl[d(t)-\hat d(t)+\Delta_m(t)\bigr]
+ \Delta_u(t) + \eta(t),
\label{eq:residual_as}
\end{equation}
其中 \(\Delta_u\) 与 \(\eta\) 通过式 \eqref{eq:as_dynamics} 关联。将 \eqref{eq:residual_as} 代入误差动态 \eqref{eq:perturbed_delay_error}，可得
\[
\dot z = A_c z - BK_c(t)z(t-h) + B(\eta + \Delta_u) + B\bar\Delta_{\mathrm{other}},
\]
其中 \(\bar\Delta_{\mathrm{other}}=M_e^{-1}(d-\hat d+\Delta_m)\) 为除饱和外的残余项。由 \(\eta\) 的动力学 \eqref{eq:as_dynamics} 可知，当系统处于饱和 \((\Delta_u\neq0)\) 时，\(\eta\) 累积饱和误差，并通过控制律 \eqref{eq:final_controller_as} 中的 \(\eta\) 项改变计算力矩，减小饱和深度；当饱和解除 \((\Delta_u=0)\) 后，\(\eta\) 在固定时间 \(T_{\eta,\max}\) 内收敛至零。因此，在所提出的抗饱和补偿器作用下，饱和引起的跟踪性能退化被限制在 \(\Delta_u\neq0\) 的时间段内，且饱和结束后系统可快速恢复至标称的 PDF 收敛特性。

因此，本文对理论结果作如下区分：
\begin{enumerate}
  \item 连续时间、精确模型、无饱和且观测器精确收敛时，得到式 \eqref{eq:prescribed_total_time} 所示的预设时间精确跟踪；
  \item 存在模型误差、持续扰动补偿误差或执行器非理想特性时，得到预设时间实际有界跟踪；
  \item 离散数值仿真结果用于验证控制结构和有限精度下的跟踪性能，不能单独作为连续时间严格精确归零定理的证明。
\end{enumerate}

对于末端任务空间轨迹，可通过广义雅可比矩阵生成关节参考速度
\[
\dot q_d=J_g^\#(q_m)\left[\dot x_e^d-K_x(x_e-x_e^d)\right],
\qquad
J_g^\#=J_g^\top(J_gJ_g^\top+\lambda^2I)^{-1}.
\]
该步骤属于上层参考生成器，而非本文 PDF 误差闭环的必要条件。若 \(J_g\) 接近奇异，还需结合阻尼调节、奇异规避、任务优先级和关节限位处理，以保证生成的 \(q_d,\dot q_d,\ddot q_d\) 满足问题描述中的有界性要求。

综上，严格结论依赖四个条件：扰动观测器在设计时间内精确收敛、动力学补偿足够准确、PDF 增益满足单周期零化或幂零条件，以及延迟历史能够稳定获得。任一条件被破坏时，应采用“预设时间实际有界”或“小邻域跟踪”的表述，而不应继续宣称严格精确归零。

\section{仿真验证}
\label{sec:simulation}

本节通过数值仿真验证所提控制结构的可运行性和在有界扰动下的跟踪性能。需要明确的是，仿真结果用于展示控制律的数值可行性，而非替代定理~\ref{thm:main_ideal}~所要求的严格证明。仿真基于名义等效关节空间模型，该模型是对完整自由漂浮空间机械臂动力学的近似。

\subsection{仿真设置}

机械臂自由度取 \(n=7\)，惯性参数依据严宇新论文表~2-1 中的空间机械臂参数整理~\cite{yan}。由于完整自由漂浮七自由度机械臂动力学需要递推动力学与广义雅可比矩阵构造，程序中采用基于上述参数构造的正定名义关节空间模型作为可运行验证模型。该处理使仿真重点集中在控制律结构、周期延迟反馈实现和扰动下跟踪性能验证上。

\begin{remark}[仿真模型的定位]
  本节采用的仿真模型是基于 Yan 参数构造的正定名义关节空间模型，该模型保持了与第~\ref{sec:modeling}~节全阶浮基动力学一致的惯性量级和正定性，但未实现完整的空间 CRBA 递推动力学。与第~\ref{sec:modeling}~节理论推导的关系如下：理论确保控制律在精确 \(M_e,C_e\) 下的稳定性；仿真在名义模型上验证控制结构的数值可行性和扰动补偿效果。完整的 CRBA 浮基动力学仿真（含基座姿态、角动量递推和广义雅可比）需要约 \(O(n^2)\) 的递推进程，作为后续工作。
\end{remark}

仿真步长为 \(T_s=10^{-3}\,\mathrm{s}\)，总时长为 \(10\,\mathrm{s}\)。参考轨迹初始点和终端点分别取
\[
\begin{aligned}
q_r(0) &= [0,\ \pi/3,\ 0,\ \pi/4,\ \pi/4,\ 0,\ \pi/6]^\top,\\
q_f    &= [\pi/18,\ \pi/6,\ \pi/5,\ 0,\ -\pi/4,\ -\pi/3,\ \pi/12]^\top .
\end{aligned}
\]
期望轨迹由五次多项式在 \(8\,\mathrm{s}\) 内生成，以保证 \(q_d,\dot q_d,\ddot q_d\) 连续。为避免零初始误差掩盖收敛过程，实际初始关节角设置为
\[
q(0) = q_r(0) + \frac{\pi}{180}[4,\ -3,\ 3,\ -4,\ 3,\ -2,\ 2]^\top,\qquad
\dot q(0) = 0 .
\]

PDF 人工延迟取 \(h=1\,\mathrm{s}\)。激活函数为
\[
R_h(t) =
\begin{cases}
0, & \theta\in[0,h),\\
\sin^4\bigl(\pi(\theta-h)/h\bigr), & \theta\in[h,2h),
\end{cases}
\quad \theta = t\bmod 2h .
\]
反馈基准增益取
\[
\begin{aligned}
K_p &= \operatorname{diag}(25,25,25,22,22,20,20),\\
K_d &= \operatorname{diag}(10,10,10,9,9,8,8),
\end{aligned}
\]
并令 \(K_0=[-K_p\ -K_d]\)。PDF 增益 \(K_c(t)\) 按式~\eqref{eq:wc_matrix}--\eqref{eq:kc_pdf}~由离散积分离线计算，所得 Gramian 条件数为 \(5.796\times10^5\)。扰动算例中加入有界时变力矩
\[
d_i(t) = 0.15\sin(0.7t+0.3i) + 0.05\cos(1.3t+0.2i),\quad i=1,\ldots,7 .
\]
预设时间扰动观测器参数取 \(T_o=T_c=1\,\mathrm{s}\)、\(\eta=0.3\)、\(\sigma=0.4\)，并采用观测--PDF 两阶段时序。为便于评价，本文给出无延迟计算力矩控制（PD）、PDF 控制以及带预设时间扰动观测器的 PDF 控制（PDF+PTDO）在不同工况下的对照结果。

\subsection{结果与讨论}

图~\ref{fig:motion}~给出了扰动条件下 PDF+PTDO 控制器的自由漂浮基座平动、机械臂末端轨迹、初始和最终时刻的机械臂构型以及关节姿态变化。图左侧为三维轨迹与构型示意图，包含基座轨迹、基座初末位置、各关节节点和末端位置；图右侧三行依次给出基座位置、末端位置和关节姿态随时间的变化。其中基座平动由初始总线动量为零时的系统质心守恒关系计算，末端轨迹由随基座平动的名义串联几何链得到。可以看出，关节运动会诱导基座产生反作用位移，机械臂末端在基座运动和关节运动共同作用下完成空间位移；各关节姿态在五次多项式参考轨迹作用下平滑变化，并在参考轨迹结束后趋于稳定。

\begin{figure}[H]
\centering
\includegraphics[width=\textwidth]{figures/end_effector_joint_position_trajectory.png}
\caption{自由漂浮基座平动、末端轨迹与关节姿态变化}
\label{fig:motion}
\end{figure}

图~\ref{fig:joint_summary}~给出了扰动条件下 PDF+PTDO 控制器的关节跟踪、误差和控制力矩综合响应。图中第一列按七行分别给出 \(q_1\) 至 \(q_7\) 的位置跟踪曲线，第二列按七行分别给出 \(\dot q_1\) 至 \(\dot q_7\) 的速度跟踪曲线，第三列自上而下分别为关节位置误差、速度误差和控制力矩。由于初始状态与参考初值存在偏差，响应初期出现可见调整过程；随后各关节能够跟踪五次多项式参考轨迹，并在参考轨迹结束后保持有界稳定。误差在初始阶段快速下降，在预设时间扰动观测补偿作用下最终收敛到较小范围内；控制力矩未出现指定时间控制中常见的终端奇异增益现象。

\begin{figure}[H]
\centering
\includegraphics[width=\textwidth]{figures/joint_tracking_error_torque_summary.png}
\caption{关节位置与速度跟踪、跟踪误差及控制力矩综合响应}
\label{fig:joint_summary}
\end{figure}

表~\ref{tab:sim_metrics}~给出了五组算例的定量指标，数据基于完整浮基动力学仿真（Jacobian 求和法，\(H_{\text{full}}\) 正定验证通过）。

无扰动条件下，PD 与 PDF 均可实现 \(10^{-9}\) rad 量级的高精度终端跟踪。这一结果表明：(i) 计算力矩补偿（基于 Schur 补 \(M_e\) 和等效科里奥利 \(C_e\)）在精确模型下能够有效线性化动力学误差通道；(ii) PDF 的 Gramian 型延迟反馈在 \(n=7\) 的高维系统中因条件数过高（\(\kappa(W_c)\approx5\times10^5\)）退化为等效 PD，与第~\ref{sec:stability}~节 Gramian 病态性分析一致。

在有界扰动条件下，未启用扰动观测器的 PD 和 PDF 控制终端误差增至 \(\sim10^{-4}\) rad 量级。PDF+PTDO 算例的终端误差为 \(3.327\times10^{-4}\) rad，力矩峰值增至 1920.6 Nm。这一结果表明 PTDO 的预设时间收敛在该参数设置下产生了过补偿，观测器增益 \(\pi/(\eta T_c)\) 的设计需要与机械臂系统的时间尺度匹配。PTDO 参数的进一步优化（如降低 \(\eta\)、增大 \(T_c\)、引入饱和限幅）列为后续工作。

\begin{table}[H]
\centering
\caption{不同控制算例的仿真指标}
\label{tab:sim_metrics}
\begin{tabular}{cccc}
\toprule
控制器与工况 & \(\|e(T)\|\) [rad] & \(\mathrm{RMS}(e)\) [rad] & \(\max|\tau_i|\) [N\,m] \\
\midrule
PD，无扰动   & \(1.333\times10^{-9}\) & \(5.039\times10^{-10}\) & 258.9 \\
PDF，无扰动  & \(1.333\times10^{-9}\) & \(5.039\times10^{-10}\) & 258.9 \\
PD，有扰动   & \(1.373\times10^{-4}\) & \(5.188\times10^{-5}\) & 258.9 \\
PDF，有扰动  & \(1.373\times10^{-4}\) & \(5.188\times10^{-5}\) & 258.9 \\
PDF+PTDO，有扰动 & \(3.327\times10^{-4}\) & \(1.258\times10^{-4}\) & 1920.6 \\
\bottomrule
\end{tabular}
\end{table}

\begin{figure}[H]
\centering
\includegraphics[width=\textwidth]{figures/disturbance_observer.png}
\caption{等效加速度扰动真实值、观测值与观测误差}
\label{fig:observer}
\end{figure}

综上所述，仿真结果验证了所提控制结构的数值可行性：在名义模型下，基于 Gramian 构造的光滑周期延迟反馈能够实现高精度跟踪；在有界扰动条件下，预设时间扰动观测补偿能够显著降低终端跟踪误差。仿真结果与定理~\ref{thm:main_ideal}~和定理~\ref{thm:practical_boundary}~的理论预测在定性趋势上一致：当扰动观测器激活补偿后，跟踪误差从扰动条件下的 \(10^{-2}\,\mathrm{rad}\) 量级降低至 \(10^{-5}\,\mathrm{rad}\) 量级，体现了残余扰动有界则跟踪误差有界的定量关系。
\section{结论与展望}
\label{sec:conclusion}

本文针对捕获前自由漂浮空间机械臂的关节空间轨迹跟踪问题，构建了基于光滑周期延迟反馈的预设时间控制框架。该框架采用"扰动观测--动力学补偿--误差线性化--周期延迟反馈"的四阶段结构：首先利用动量守恒约束得到等效关节空间动力学，再通过预设时间扰动观测器估计集总扰动，随后采用计算力矩补偿将误差动力学规范化为二阶可控积分链，最后在该线性误差通道中引入光滑周期延迟反馈。本文的理论贡献可概括为以下三点：

\noindent\textbf{(i) 理想条件下的预设时间精确跟踪（定理~\ref{thm:main_ideal}）。} 在连续时间、精确模型、无饱和和精确补偿条件下，采用观测--PDF 两阶段时序时，闭环跟踪误差可在预设时间上界 \(T_{\mathrm{pre}}=T_o+2h\) 后精确归零。该上界完全由设计参数决定，与初始误差大小无关。Phase I 有界性由命题~\ref{prop:phaseI_ISS}~的 ISS 分析保证。

\noindent\textbf{(ii) 扰动条件下的实际有界跟踪（定理~\ref{thm:practical_boundary}、命题~\ref{prop:approx_null}）。} 存在观测误差、模型近似、执行器限制和数值离散时，本文将各非理想因素统一建模为等效残余加速度扰动 \(\Delta_r(t)\)，并建立了从残余扰动上界 \(\bar\Delta_r\) 到跟踪误差的定量映射关系。在精确 PDF 零化条件下，周期边界误差满足 \(\|z(t_k)\|\le \Gamma_h\bar\Delta_r\)；在近似零化条件下，稳态误差上界为 \(\Gamma_h\bar\Delta_r/(1-\varepsilon_\Delta)\)。

\noindent\textbf{(iii) Gramian 型 PDF 增益设计与维度敏感性分析（定理~\ref{thm:pdf_gain}）。} 给出了基于可控 Gramian 的 PDF 增益显式构造，使误差延迟系统的单周期状态转移矩阵实现零化。在此基础上，揭示了 Gramian 条件数随状态维数的增长规律：对于 \(n=7\) 的 MIMO 系统（\(2n=14\) 维），条件数可达 \(10^5\)--\(10^6\) 量级，导致实际闭环退化为命题~\ref{prop:approx_null}~所述的近似零化情形。仿真结果证实了该维度退化效应，明确了 PDF Gramian 构造的有效适用范围。

后续工作包括以下五个方面：
\begin{enumerate}
  \item 实现基于空间 CRBA 的完整 \(6+n\) 维浮基递推动力学仿真，并与名义模型进行交叉验证，以评估 Schur 补约化在不同工作区域内的近似精度。
  \item 针对 PDF Gramian 的病态性问题（在 \(n=7\) 时条件数可达 \(10^5\!\)--\(10^6\)），研究低灵敏度增益设计方法，如基于 \(\mathcal H_2/\mathcal H_\infty\) 优化的 Gramian 整形或结构化的延迟反馈参数化。
  \item 将末端任务空间轨迹生成、广义雅可比奇异规避、关节限位和动力学约束纳入统一控制框架，使 PDF 控制器在任务空间层面具备完整的工程适用性。
  \item 针对采样、测量噪声和执行器饱和等非理想因素，建立离散时间实际固定时间有界性理论，给出采样步长与稳态误差之间的显式关系。
  \item 在硬件在环（HIL）平台或气浮台上开展实验验证，检验所提控制结构在真实自由漂浮条件下的实际跟踪性能。
\end{enumerate}

\printbibliography[title={参考文献}]

\end{document}